\section{Definizione del Problema}

\subsection{La complessità intrinseca delle NullPointerException }
Le \textit{NullPointerException} (NPE) rappresentano una delle cause principali di eccezioni
 non gestite che determinano il crash delle applicazioni in ambienti di produzione 
 \cite{durieux2017dynamic}. In ottica ingegneristica, è fondamentale chiarire che le 
 NPE non costituiscono banali errori sintattici facilmente individuabili, bensì insidiosi 
 bug semantici legati alla gestione della memoria e dei puntatori. La loro insorgenza e il
  loro comportamento sono per loro natura legati e dipendenti dal flusso di esecuzione 
  (control-flow) e dal flusso dei valori (value-flow) del programma \cite{xu2019vfix}. 
  Come evidenziato dalla letteratura, l'analisi statica del codice locale si rivela spesso 
  insufficiente per avere un quadro completo degli oggetti e delle variabili compatibili;
   per comprendere appieno la natura di una NPE risulta infatti necessaria un'analisi dinamica 
   e dettagliata del contesto a tempo di esecuzione \cite{durieux2017dynamic}.

La difficoltà principale nel debugging delle NPE risiede in un marcato disallineamento spaziale 
e logico tra il difetto originario e il punto di interruzione in cui l'eccezione si 
manifesta. Il punto esatto in cui il programma va materialmente in crash
rappresenta tipicamente solo la fine di una complessa catena di propagazione originata altrove. 
Un difetto originario, che nella maggior parte dei casi consiste in una mancata inizializzazione 
di un oggetto o in un controllo di nullità assente, risiede spesso a notevole distanza a livello 
di codice \cite{xu2019vfix}. Di conseguenza, una NPE e la sua effettiva posizione di correzione 
(bug-fixing location) possono estendersi in modo inter-procedurale attraverso molteplici funzioni 
all'interno di codebase di grandi dimensioni \cite{xu2019vfix}.

Questa frammentazione strutturale pone sfide profonde per gli strumenti di
 \textit{Automated Program Repair}. Ignorando le complesse dipendenze 
 di controllo e dei dati (il value-flow inter-procedurale), i tradizionali 
 approcci APR operano su spazi di ricerca immensi generando spesso riparazioni 
 localizzate esclusivamente al punto di crash \cite{xu2019vfix}. Tali approcci 
 limitati (come i \textit{single-line edits} per saltare una singola istruzione 
 o inserire un controllo banale) portano frequentemente alla generazione di patch 
 "plausibili ma errate" (\textit{plausible but incorrect patches}): soluzioni che 
 riescono a superare la suite di test ma introducono comportamenti semanticamente 
 indefiniti nel prosieguo del programma \cite{xu2019vfix}. Risolvere in modo definitivo
  questi bug complessi non richiede un semplice mascheramento dell'eccezione, ma necessita 
  di modifiche strutturate (\textit{multi-point edits}) capaci di ripristinare il 
  corretto flusso dei dati fin dalla sua origine, evidenziando così i limiti degli 
  strumenti di riparazione basati unicamente su template sintattici limitati o approcci 
  di ricerca ciechi \cite{xu2019vfix}.

\subsection{I limiti degli LLM nell'Automated Program Repair}

Negli ultimi anni, i \textit{Large Language Models} si sono affermati come lo standard de facto 
per automatizzare \textit{task} complessi di ingegneria del software, incluse le più recenti tecniche 
di \textit{Automated Program Repair} a livello di \textit{repository} \cite{zhang2024llmhallucinations,pabba2025semagent}.
Tuttavia, nonostante le notevoli capacità dimostrate nella risoluzione di frammenti di codice 
isolati, un'analisi accademica rigorosa rivela che gli LLM "puri" (intendendo con ciò l'utilizzo 
dei \textit{Large Language Models} in totale isolamento, affidandosi esclusivamente alla loro architettura 
interna e alle informazioni acquisite durante la fase di addestramento) presentano dei difetti strutturali 
intrinseci quando sono chiamati ad affrontare \textit{bug} complessi e semanticamente densi. Tali limiti
emergono in modo evidente di fronte a problematiche che richiedono una profonda comprensione del 
flusso di controllo, dei dati o delle dipendenze di stato, come ad esempio \textit{bug} critici legati a 
\textit{NullPointerException} o altre violazioni della logica di esecuzione profonda \cite{zhang2025semanticforge}.

Il motivo tecnico alla base di questi difetti strutturali risiede nel meccanismo fondamentale di
 funzionamento dei modelli linguistici. Come evidenziato da Zhang et al., gli LLM generano il 
 codice operando primariamente sulla base della probabilità statistica, ovvero prevedendo ricorsivamente 
 un \textit{token} dopo l'altro. 
 Questo processo avviene senza alcun meccanismo connaturale capace di garantire il rispetto dei requisiti 
formali o semantici del linguaggio di programmazione \cite{zhang2025semanticforge}.
Precisano, inoltre, che il vantaggio principale degli LLM, ovvero l'eccellente capacità 
di riconoscimento dei \textit{pattern}, costituisce al contempo la loro più grande limitazione: 
i modelli tendono a generare codice appoggiandosi esclusivamente ai \textit{pattern} comuni osservati 
nei dati di addestramento, piuttosto che derivare una comprensione profonda delle specifiche 
intenzioni e necessità logiche del problema \cite{zhang2024llmhallucinations}.
Questo approccio puramente predittivo conduce inesorabilmente al fenomeno delle allucinazioni \cite{zhang2025semanticforge,zhang2024llmhallucinations}.
In ambito di generazione del codice, le allucinazioni si manifestano tipicamente quando il modello 
produce \textit{patch} che appaiono sintatticamente corrette e del tutto plausibili, spesso riuscendo 
persino a superare la fase di compilazione senza errori, ma che si rivelano fondamentalmente 
errate o inconsistenti dal punto di vista della logica semantica. Tali aberrazioni si dividono 
in errori logici (come l'iterazione su collezioni errate o l'omissione di mutazioni di stato vitali) 
ed errori schematici (che includono violazioni di firme, discrepanze di tipo e chiamate a funzioni 
mai definite) \cite{zhang2025semanticforge}.

Questo limite statistico e predittivo si traduce direttamente nel fenomeno dell'\textit{overfitting}, 
ampiamente documentato da Pabba et al. Gli studi dimostrano come i sistemi basati esclusivamente 
su LLM tendano a iper-localizzarsi sulle righe di codice immediatamente sospette, tentando 
di ripararle in totale isolamento rispetto all'intero ecosistema del software. Il risultato 
è la generazione di \textit{patch} che soffrono di un marcato \textit{overfitting}: il codice prodotto riesce 
a risolvere gli specifici esempi falliti forniti nella descrizione del \textit{bug}, riuscendo magari 
a superare la \textit{test suite} locale, ma si rivela incapace di generalizzare. La riparazione 
si limita quindi ad agire come un mero cerotto superficiale che fallisce nel gestire i casi 
limite (\textit{edge cases}) e non soddisfa i requisiti di coerenza più ampi e profondi dell'applicazione \cite{pabba2025semagent}.

Infine, questa tendenza all'\textit{overfitting} e alla generazione di soluzioni localmente ottime ma 
globalmente fallate è inestricabilmente legata alla mancanza di una reale consapevolezza 
del codice. Zhang et al. sottolineano come l'assenza di un contesto a livello di intero progetto 
(\textit{repository-level context}) impedisca agli LLM puri di comprendere le dipendenze profonde necessarie 
per vere riparazioni \cite{zhang2024llmhallucinations}. In scenari di sviluppo reali, la stragrande maggioranza delle funzioni 
dipende da entità definite altrove nel progetto o da API personalizzate; tuttavia, i rigidi limiti 
sui \textit{token} dei modelli \textit{Transformer} rendono impossibile e impraticabile l'inclusione dell'intero 
contesto del \textit{repository} nel \textit{prompt}. Di conseguenza, ignorando le catene di chiamate transitive 
e i vincoli architetturali distribuiti su molteplici file, i modelli "puri" si affidano 
a supposizioni statistiche cieche, generando conflitti di dipendenza e riparazioni semanticamente 
vuote che non possono integrarsi in sicurezza in basi di codice complesse \cite{zhang2024llmhallucinations,zhang2025semanticforge}.

\subsection{Il vizio di forma della ricerca: l'illusione della Perfect Fault Localization}

\subsection{La necessità di una guida semantica rigorosa}
